\begin{frame}[shrink]
	\frametitle{tdBN: colocar a corrente na escala do limiar}
	\small
	\par BatchNorm comum padroniza valores para média 0 e variância 1. Isso funciona em redes contínuas, mas um neurônio de pulso só dispara quando seu potencial alcança $V_{th}$; uma escala fixa pode deixar toda a camada abaixo desse limiar.
	\begin{center}
		\begin{tabular}{lcc}
			\toprule
			& \textbf{BatchNorm comum} & \textbf{tdBN, $V_{th}=2$, $\alpha=1$} \\
			\midrule
			Entrada $X=[0,2,4,6]$ & $[-1{,}34,-0{,}45,0{,}45,1{,}34]$ & $[-2{,}68,-0{,}89,0{,}89,2{,}68]$ \\
			Escala do desvio & $1$ & $\alpha V_{th}=2$ \\
			Valores acima do limiar & $0$ de $4$ & $1$ de $4$ \\
			\bottomrule
		\end{tabular}
	\end{center}
	\begin{itemize}
		\item A tdBN usa $Y=\alpha V_{th}(X-\mu)/\sigma$: a escala acompanha o limiar, em vez de ignorá-lo.
		\item \textbf{Por que importa?} Mais valores podem alcançar o limiar e gerar gradientes úteis; isso ajuda a treinar redes de pulso profundas e reduz o problema de camadas sem disparos.
		\item Não é garantia de que cada neurônio disparará nem substitui uma boa escolha de dinâmica e perda. É uma normalização proposta para estabilizar o treinamento direto de SNNs profundas \cite{zheng2021deeper, eshraghian2023training}.
	\end{itemize}
\end{frame}
